% Lab 3 of Intelligent Robotics (ASY2): a linear Kalman filter that finds a driving robot from GPS and odometry.
% Build: make  (latexmk -pdf). Figures are TikZ files in figures/, made by tools/make_figures.py.
\documentclass[11pt,a4paper,parskip=half]{scrartcl}
\usepackage[english]{babel}
\usepackage[T1]{fontenc}
\usepackage[utf8]{inputenc}
\usepackage[sfdefault,lining]{FiraSans}
\usepackage[scale=0.88]{FiraMono}
\usepackage[a4paper,left=2.4cm,right=2.4cm,top=2.6cm,bottom=2.6cm,headsep=0.8cm]{geometry}
\usepackage{amsmath,amssymb}
\usepackage{booktabs,array,tabularx}
\usepackage{xcolor}
\usepackage{listings}
\usepackage{tikz}
\usetikzlibrary{arrows.meta,calc,shapes.geometric}
\usepackage[skins,breakable]{tcolorbox}
\usepackage{fontawesome5}
\usepackage[font=small,labelfont={bf,color=ohmred},skip=4pt]{caption}
\usepackage{enumitem}
\usepackage{microtype}
\usepackage[headsepline=0.4pt]{scrlayer-scrpage}
\usepackage[colorlinks=true,linkcolor=ohmred,urlcolor=ohmred]{hyperref}
\usepackage{cleveref}

% ---------------------------------------------------------------- colours and headings
\definecolor{ohmred}{HTML}{C8372D}
\definecolor{ohmgreen}{HTML}{2E8B57}
\definecolor{ohmamber}{HTML}{D98E04}
\definecolor{ohmgrey}{HTML}{5A5A5A}
\addtokomafont{disposition}{\color{ohmred}}
\addtokomafont{section}{\Large}
\addtokomafont{subsection}{\large}
\setkomafont{pageheadfoot}{\small\color{ohmgrey}}
\setkomafont{pagenumber}{\small\bfseries\color{ohmred}}
\clearpairofpagestyles
\ihead{Intelligent Robotics (ASY2) · Lab 3}
\ohead{Kalman filter}
\ofoot{\pagemark}
\pagestyle{scrheadings}

\lstset{basicstyle=\ttfamily\small, columns=fullflexible, keepspaces=true,
        showstringspaces=false, breaklines=true, frame=none, backgroundcolor=\color{black!5},
        xleftmargin=8pt, xrightmargin=8pt, framexleftmargin=6pt, framexrightmargin=6pt,
        aboveskip=4pt, belowskip=4pt}

\newcommand{\thema}[1]{\texttt{#1}}
\newcommand{\datei}[1]{\texttt{#1}}
\newcommand{\minutes}[1]{\hfill{\normalsize\mdseries\color{ohmgrey}\faClock[regular]~#1 min}}
\newcommand{\mat}[1]{\mathbf{#1}}
\newcommand{\T}{^{\mathsf{T}}}

% ---------------------------------------------------------------- boxes
\tcbset{base/.style={enhanced, breakable, boxrule=0pt, arc=1.5pt, left=8pt, right=8pt,
        top=5pt, bottom=5pt, before skip=8pt, after skip=8pt}}
% \begin{todo}{1.3}: the block STUDENT 1.3, with its hint command at the end of the last line
\newcommand{\hintcmd}[1]{\unskip\nobreak\hfil\penalty50\hskip1em\null\nobreak\hfil
        \mbox{\small\color{ohmgrey}\faLightbulb[regular]~\texttt{python3 hint.py #1}}{\parfillskip=0pt\par}}
\newtcolorbox{todo}[1]{base, colback=ohmred!5, borderline west={2.5pt}{0pt}{ohmred},
        title={\faPen~STUDENT #1}, coltitle=ohmred, colbacktitle=ohmred!5, fonttitle=\bfseries,
        attach title to upper={\quad}, after upper={\hintcmd{#1}}}
\newtcolorbox{script}[1]{base, colback=black!3, borderline west={2.5pt}{0pt}{ohmgrey},
        title={\faBook~Lecture script, #1}, coltitle=ohmgrey, colbacktitle=black!3,
        fonttitle=\small\bfseries, toptitle=2pt, bottomtitle=0pt}
\newtcolorbox{note}{base, colback=ohmamber!8, borderline west={2.5pt}{0pt}{ohmamber}}
\newtcolorbox{checkpoint}{base, colback=ohmgreen!7, borderline west={2.5pt}{0pt}{ohmgreen},
        title={\faCheckCircle[regular]~Checkpoint}, coltitle=ohmgreen, colbacktitle=ohmgreen!7,
        fonttitle=\bfseries, attach title to upper={\quad}}

% \running{<launch arguments>}{<node arguments>}: what terminals 1 and 2 run during a task
\newcommand{\running}[2]{\begin{tcolorbox}[base, colback=black!3, borderline west={2.5pt}{0pt}{ohmgrey},
        title={\faTerminal~Running}, coltitle=ohmgrey, colbacktitle=black!3, fonttitle=\small\bfseries,
        toptitle=2pt, bottomtitle=0pt, fontupper=\small]
  \begin{tabularx}{\linewidth}{@{}>{\color{ohmgrey}}l >{\raggedright\arraybackslash}X@{}}
    terminal 1 & \texttt{ros2 launch ohm\_kalman\_localization kalman.launch.py#1} \\
    terminal 2 & \texttt{ros2 run ohm\_kalman\_localization kalman\_node.py -\kern0pt-ros-args #2}
  \end{tabularx}
\end{tcolorbox}}
% \expected[<note>]{<name>}: the path (figures/run_<name>.tikz) and the monitor (figures/term_<name>.txt)
\newtcolorbox{seebox}{base, colback=ohmgreen!3, borderline west={2.5pt}{0pt}{ohmgreen!60},
        title={\faEye[regular]~What you should see}, coltitle=ohmgreen!80!black, colbacktitle=ohmgreen!3,
        fonttitle=\small\bfseries, toptitle=2pt, bottomtitle=0pt, breakable=false}
\newcommand{\expected}[2][]{\begin{seebox}
  \centering\resizebox{0.5\linewidth}{!}{\input{figures/run_#2.tikz}}\par\smallskip
  \raggedright{\footnotesize\color{ohmgrey}#1the last 70\,s: white the truth, green the fixes, red your
    estimate and its \(1\sigma\) ellipse. Terminal 1, the monitor (the numbers change from lap to lap):}
  \lstinputlisting[basicstyle=\ttfamily\scriptsize, aboveskip=2pt, belowskip=0pt, xleftmargin=4pt,
                   framexleftmargin=3pt, breakindent=1.5em]{figures/term_#2.txt}
\end{seebox}}

\setlist[enumerate,1]{leftmargin=*,itemsep=3pt,topsep=2pt}
\setlist[itemize,1]{leftmargin=*,itemsep=2pt,topsep=2pt}

\begin{document}
\sloppy

% ---------------------------------------------------------------- title
\begin{tcolorbox}[enhanced, colback=ohmred, colframe=ohmred, arc=0pt, left=14pt, right=14pt,
                  top=12pt, bottom=12pt, fontupper=\color{white}]
  {\small Intelligent Robotics (ASY2) · Lab 3}\\[6pt]
  {\LARGE\bfseries Kalman filter}\\[4pt]
  {\large Where is the robot? GPS and odometry, fused}
\end{tcolorbox}

\noindent
\begin{tabular}{@{}>{\color{ohmgrey}}l l@{\hspace{2.4em}}>{\color{ohmgrey}}l l@{}}
  \faClock[regular]~Duration & 180 minutes & \faUsers~Work in & pairs \\
  \faBook~Script & Cycle 3, section 3.1 & \faFolderOpen[regular]~Files & \datei{ohm-kalman-localization/student/}
\end{tabular}

\section{Overview}
A robot drives a loop on an open field, lap after lap (\cref{fig:field}). You do not steer it: a driver node
does. Your node has to find out where the robot is. It has two sensors:
\begin{itemize}
  \item \textbf{GPS}: five fixes per second. Each fix scatters by \(\sigma = 0.5\)\,m in \(x\) and in \(y\).
        It does not drift. In the tunnel there is no fix.
  \item \textbf{Odometry}: 50 messages per second from the wheels. Its velocity is good. Its position drifts,
        because the wheel constants are a little wrong.
\end{itemize}
You write a Kalman filter that combines a motion model with these measurements:
\begin{itemize}
  \item \textbf{Task 1}: GPS alone, with a constant-velocity model.
  \item \textbf{Task 2}: GPS and the velocity of the odometry. Then the tunnel.
\end{itemize}

\paragraph{Why a linear filter is enough.} The state is \(\vec x = (x, y, v_x, v_y)\T\) in the world frame.
The constant-velocity model is linear. The GPS measures \(x\) and \(y\): linear. The odometry measures the
velocity in the frame of the robot; rotated into the world frame with the heading of the odometry, it is
\(v_x\) and \(v_y\): linear again. So the plain Kalman filter of the script applies. The heading itself is not
estimated, it is taken from the odometry. Estimating it as well makes the model non-linear: that is the EKF
(\cref{sec:outlook}).

\begin{figure}[ht]
  \centering
  % generated by tools/make_figures.py: the hall open, 1 m = 0.4 cm
\begin{tikzpicture}[x=0.4cm, y=0.4cm, font=\footnotesize]
  \fill[black!4] (0,0) rectangle (30,20);
  \draw[black!50, line width=1.2pt] (0,0) rectangle (30,20);
  \draw[black!12] (0,0) grid[step=5] (30,20);
  \fill[fill=black!35] (7.5,9.0) rectangle (12.0,14.5);
  \node[above, align=center] at (9.75,14.5) {tunnel};
  \draw[ohmred, line width=1pt] (12.00,7.00) -- (12.25,7.00) -- (12.50,7.00) -- (12.75,7.00) -- (13.00,7.00) -- (13.25,7.00) -- (13.50,7.00) -- (13.75,7.00) -- (14.00,7.00) -- (14.25,7.00) -- (14.50,7.00) -- (14.75,7.00) -- (15.00,7.00) -- (15.25,7.00) -- (15.50,7.00) -- (15.75,7.00) -- (16.00,7.00) -- (16.25,7.00) -- (16.50,7.00) -- (16.75,7.00) -- (17.00,7.00) -- (17.25,7.00) -- (17.50,7.00) -- (17.75,7.00) -- (18.00,7.00) -- (18.25,7.00) -- (18.50,7.00) -- (18.75,7.00) -- (19.00,7.00) -- (19.25,7.00) -- (19.50,7.00) -- (19.75,7.00) -- (20.00,7.00) -- (20.26,7.03) -- (20.50,7.13) -- (20.71,7.29) -- (20.87,7.50) -- (20.97,7.74) -- (21.00,8.00) -- (21.00,8.25) -- (21.00,8.50) -- (21.00,8.75) -- (21.00,9.00) -- (21.00,9.25) -- (21.00,9.50) -- (21.00,9.75) -- (21.00,10.00) -- (21.00,10.25) -- (21.00,10.50) -- (21.00,10.75) -- (21.00,11.00) -- (21.00,11.25) -- (21.00,11.50) -- (21.00,11.75) -- (21.00,12.00) -- (20.97,12.26) -- (20.87,12.50) -- (20.71,12.71) -- (20.50,12.87) -- (20.26,12.97) -- (20.00,13.00) -- (19.75,13.00) -- (19.50,13.00) -- (19.25,13.00) -- (19.00,13.00) -- (18.75,13.00) -- (18.50,13.00) -- (18.25,13.00) -- (18.00,13.00) -- (17.75,13.00) -- (17.50,13.00) -- (17.25,13.00) -- (17.00,13.00) -- (16.75,13.00) -- (16.50,13.00) -- (16.25,13.00) -- (16.00,13.00) -- (15.75,13.00) -- (15.50,13.00) -- (15.25,13.00) -- (15.00,13.00) -- (14.75,13.00) -- (14.50,13.00) -- (14.25,13.00) -- (14.00,13.00) -- (13.75,13.00) -- (13.50,13.00) -- (13.25,13.00) -- (13.00,13.00) -- (12.75,13.00) -- (12.50,13.00) -- (12.25,13.00) -- (12.00,13.00) -- (11.75,13.00) -- (11.50,13.00) -- (11.25,13.00) -- (11.00,13.00) -- (10.75,13.00) -- (10.50,13.00) -- (10.25,13.00) -- (10.00,13.00) -- (9.74,12.97) -- (9.50,12.87) -- (9.29,12.71) -- (9.13,12.50) -- (9.03,12.26) -- (9.00,12.00) -- (9.00,11.75) -- (9.00,11.50) -- (9.00,11.25) -- (9.00,11.00) -- (9.00,10.75) -- (9.00,10.50) -- (9.00,10.25) -- (9.00,10.00) -- (9.00,9.75) -- (9.00,9.50) -- (9.00,9.25) -- (9.00,9.00) -- (9.00,8.75) -- (9.00,8.50) -- (9.00,8.25) -- (9.00,8.00) -- (9.03,7.74) -- (9.13,7.50) -- (9.29,7.29) -- (9.50,7.13) -- (9.74,7.03) -- (10.00,7.00) -- (10.25,7.00) -- (10.50,7.00) -- (10.75,7.00) -- (11.00,7.00) -- (11.25,7.00) -- (11.50,7.00) -- (11.75,7.00) -- cycle;
  \draw[-{Stealth[length=5pt]}, ohmred, line width=1.4pt] (12.0,7.0) -- (14.0,7.0);
  \node[below] at (12.0,7.0) {start};
  \fill[ohmred] (9.0,10.0) circle (0.35);
  \node[left] at (8.7,10.0) {stop};
  \draw[-{Stealth[length=4pt]}] (1,1) -- (4,1) node[right] {$x$};
  \draw[-{Stealth[length=4pt]}] (1,1) -- (1,4) node[above] {$y$};
  \node[below left] at (30,0) {30\,m};
  \node[above right, rotate=90] at (30,0) {};
  \node[above left] at (0,20) {20\,m};
\end{tikzpicture}

  \caption{The hall \thema{open} from above, 30\,m \(\times\) 20\,m. The robot drives the red loop
    counter-clockwise at 0.5\,m/s, about 70\,s per lap. Every lap it brakes, stands for 4\,s at the stop and
    starts again. Grey: the tunnel (\thema{gps:=tunnel}). Amber: a glass front that reflects the GPS signal
    (\thema{gps:=multipath}, extra task).}
  \label{fig:field}
\end{figure}

\begin{table}[ht]\centering
\caption{Schedule.}
\label{tab:schedule}
\begin{tabular}{@{}rll@{}}
\toprule
min & task & done when \\
\midrule
0--15 & \cref{sec:task0} The sensors & \cref{tab:sensors} is filled in \\
15--55 & \cref{sec:task1} GPS alone: the filter & \thema{pytest -k task1} passes \\
55--70 & \cref{sec:task1b} GPS alone: in the loop & the red path follows the robot \\
70--100 & \cref{sec:task1c} GPS alone: the process noise & \cref{tab:q} is filled in \\
100--110 & break & \\
110--145 & \cref{sec:task2} GPS and odometry & \cref{tab:sigmav} is filled in \\
145--170 & \cref{sec:task2b} The tunnel & \cref{tab:tunnel} is filled in \\
170--180 & \cref{sec:extra} Extra task: multipath & \\
\bottomrule
\end{tabular}
\end{table}

\paragraph{Notation} as in the script. In the code the names are short:
\begin{center}\small
\begin{tabular}{@{}lll@{}}
\toprule
script & code & \\
\midrule
\(\tilde{\vec x}_n\), \(\tilde{\vec x}_n^+\) & \thema{self.x} & state before and after the correction \\
\(\tilde{\mat P}_n\), \(\tilde{\mat P}_n^+\) & \thema{self.P} & its covariance, \(4 \times 4\) \\
\(\mat A\), \(\mat Q\) & \thema{A}, \thema{Q} & model and process noise of one step of \(\Delta t\) = \thema{dt} seconds \\
\(\hat{\vec y}_n\) & \thema{z} & the measurement \\
\(\hat{\vec y}_n - \tilde{\vec y}_n\) & \thema{innovation} & measured minus predicted measurement \\
\(\mat H\), \(\mat R\), \(\mat K_n\) & \thema{H}, \thema{R}, \thema{K} & observation matrix, measurement noise, gain \\
\bottomrule
\end{tabular}
\end{center}
Everything is in metres and seconds, in the world frame \thema{map}.

\newpage
\section{Getting started}
\paragraph{Your own copy.} Fork \datei{ohm-kalman-localization} on GitHub (button \emph{Fork} at the top
right), then clone your fork into the \datei{src/} folder of your workspace. Update the simulator:
\begin{lstlisting}
git clone git@github.com:<your-name>/ohm-kalman-localization.git
\end{lstlisting}
\begin{lstlisting}
git -C ir-mecanum-lab pull
\end{lstlisting}
Commit after every task and push, so nothing is lost at the end of the session.

\paragraph{Build.} Build once, then source the workspace in every new terminal:
\begin{lstlisting}
colcon build --symlink-install --packages-select mecanum_lab_interfaces mecanum_lab ohm_kalman_localization
\end{lstlisting}
\begin{lstlisting}
source install/setup.bash
\end{lstlisting}
You do not need to build again after a change to your code. Just start the node again.

\paragraph{Simulator.} One launch file starts the simulator in the hall \thema{open}, the driver, the
monitor and RViz. The robot starts to drive at once.
\begin{lstlisting}
ros2 launch ohm_kalman_localization kalman.launch.py
\end{lstlisting}
The driver steers by the exact pose on \thema{/muster/truth}. Your filter never reads \thema{/truth}: it
has to find the robot from GPS and odometry alone. The monitor compares your estimate with the truth.

\paragraph{Your files.} \datei{student/kalman\_filter.py} (prediction and correction) and
\datei{student/models.py} (the model and the two sensors). You write only the few lines in the middle, marked
\begin{lstlisting}
# the state through the model (no input B u here), then the covariance through the model plus Q
# TODO STUDENT 1.2
raise NotImplementedError("STUDENT 1.2")
\end{lstlisting}
The comment above it says in words what goes there; the equation is in the task of this handout. Replace the
\thema{raise} line by your code.

\begin{note}
\textbf{numpy} for the matrices, as in Matlab: \thema{np.array([[1, 2], [3, 4]])} is a matrix, row by row;
\thema{A @ B} the matrix product (\thema{A * B} multiplies element by element: not what you want);
\thema{A.T} the transpose; \thema{np.linalg.inv(S)} the inverse; \thema{np.eye(4)} the identity;
\thema{np.diag([a, b])} a diagonal matrix; \thema{P[0, 2]} row 0, column 2.
\end{note}

\paragraph{The node.} \datei{student/kalman\_node.py} is given. It subscribes to the sensors, calls your
functions and publishes the estimate. The parameter \thema{task} says which sensors it uses. It runs in a
second terminal:
\begin{lstlisting}
ros2 run ohm_kalman_localization kalman_node.py --ros-args -p task:=1
\end{lstlisting}
If a function is not written yet, the node stops with \thema{NotImplementedError: STUDENT 1.1}: that is the
next one to write.

\paragraph{Self check.} The tests try your functions on made-up numbers, one task at a time:
\begin{lstlisting}
python3 -m pytest test/test_filter.py -q -k task1
\end{lstlisting}

\paragraph{Three terminals.} All of them in the workspace, sourced.
\begin{center}\small
\begin{tabular}{@{}>{\color{ohmgrey}}l l l@{}}
  terminal 1 & simulator, driver, monitor, RViz & the monitor prints here every 10\,s \\
  terminal 2 & your node           & Ctrl-C and start again after every change \\
  terminal 3 & tests, hints, plots, \thema{ros2 topic} & \\
\end{tabular}
\end{center}
Every task starts with a box \emph{Running}: what terminals 1 and 2 run. It ends with \emph{What you
should see}: the last 70\,s of the path and the monitor when the task is done. RViz uses the same colours:
white the true path, green the GPS fixes, red your estimate with its \(1\sigma\) ellipse.

\paragraph{The monitor and the plots.} Every 10\,s the monitor prints, over the last 30\,s, the error (RMSE) of
the GPS and of your filter against the truth, the largest error, the \(\sigma\) your filter claims, and the NEES
(\cref{sec:task1c}).
Each start of your node is a new run: the monitor writes it to \datei{runs/}. The last run as plots, like the
plots of the Octave tutorial (path, velocity, diagonal of \(\mat P\) over time):
\begin{lstlisting}
ros2 run ohm_kalman_localization plot_run.py
\end{lstlisting}

\paragraph{Hints.} Stuck? Every block has three hints: the idea, how to write it in Python, and at last
the code. Each call opens the next one, in the folder \datei{ohm-kalman-localization}:
\begin{lstlisting}
python3 hint.py 1.3
\end{lstlisting}
\thema{python3 hint.py} alone shows the hints you have opened. Try on your own first: the hints notice
when you ask too early.

\begin{figure}[!hb]
  \centering
  \resizebox{\linewidth}{!}{% The ROS graph of the lab. Red: your node. Grey: given.
\begin{tikzpicture}[font=\footnotesize, >={Stealth[length=5pt]},
    node/.style={rectangle, rounded corners=2pt, draw=black!50, fill=black!4, minimum height=0.9cm,
                 minimum width=2.4cm, align=center},
    mine/.style={node, draw=ohmred, fill=ohmred!8, line width=0.9pt},
    topic/.style={font=\scriptsize\ttfamily, fill=white, inner sep=1.5pt}]
  \node[node] (sim) at (0,0) {simulator\\\scriptsize\texttt{mecanum\_sim}};
  \node[mine] (kf) at (6.2,0) {your filter\\\scriptsize\texttt{kalman\_node.py}};
  \node[node] (mon) at (11.6,0) {monitor\\\scriptsize\texttt{monitor\_node.py}};
  \node[node] (drv) at (0,-2.4) {driver\\\scriptsize\texttt{route\_node.py}};
  \node[node] (rviz) at (11.6,-2.4) {RViz};
  \draw[->] (sim.20) -- node[topic, above=1pt] {/muster/gps\_cov} (kf.160);
  \draw[->] (sim.-20) -- node[topic, below=1pt] {/muster/odom (task 2)} (kf.200);
  \draw[->] (kf.20) -- node[topic, above=1pt] {/muster/kf/pose} (mon.160);
  \draw[->] (kf.-20) -- node[topic, below=1pt] {/muster/kf/twist} (mon.200);
  \draw[->] (sim.250) -- node[topic, left=1pt] {/muster/truth} (drv.110);
  \draw[->] (drv.70) -- node[topic, right=1pt] {/muster/cmd\_vel} (sim.290);
  \draw[->, black!50] (sim.90) -- ++(0,0.7) -| node[topic, pos=0.25, above=1pt] {/muster/truth, /muster/gps\_cov}
        (mon.90);
  \draw[->] (mon) -- node[topic, right=1pt] {/muster/kalman\_lab/*} (rviz);
\end{tikzpicture}
}
  \caption{The nodes and topics of the lab. Red: your node. Grey: given.}
  \label{fig:nodes}
\end{figure}

\clearpage
\section{Task 0: the sensors\minutes{15}}
\label{sec:task0}
\running{}{(nothing yet)}

Before you filter, look at what you get. In terminal 3:
\begin{lstlisting}
ros2 topic list
\end{lstlisting}
\begin{lstlisting}
ros2 topic hz /muster/gps_cov
\end{lstlisting}
\begin{lstlisting}
ros2 topic echo --once /muster/gps_cov
\end{lstlisting}
\begin{lstlisting}
ros2 interface show geometry_msgs/msg/PoseWithCovarianceStamped
\end{lstlisting}
Do the same for \thema{/muster/odom} (type \thema{nav\_msgs/msg/Odometry}). Fill in \cref{tab:sensors}.

\begin{table}[ht]\centering
\renewcommand{\arraystretch}{1.5}
\caption{The two sensors. Fill in.}
\label{tab:sensors}
\begin{tabular}{@{}>{\raggedright\arraybackslash}p{6.0cm}p{3.6cm}p{3.6cm}@{}}
\toprule
& \thema{/muster/gps\_cov} & \thema{/muster/odom} \\
\midrule
message type & & \\
messages per second & & \\
frame of the position (\thema{header.frame\_id}) & & \\
frame of the velocity (\thema{child\_frame\_id}) & --- & \\
\(\sigma\) of the GPS position, from \thema{covariance} & & --- \\
speed while the robot stands at the stop & --- & \\
\bottomrule
\end{tabular}
\end{table}

\begin{itemize}
  \item The covariance of the GPS fix is a list of 36 numbers. Which of them is the variance of \(x\), which
        the variance of \(y\)? You need them in \cref{sec:task1b}.
  \item The simulator window shows the GPS fixes (key \thema{g}). Press \thema{o}: now it also shows the robot
        as the odometry sees it. Watch both for a lap. Which one scatters, which one drifts?
\end{itemize}

\begin{checkpoint}
\(\sigma\) of the GPS is 0.5\,m. The odometry gives the velocity in the frame of the robot, not in the world.
\end{checkpoint}

\newpage
\section{Task 1: GPS alone, the filter\minutes{40}}
\label{sec:task1}
\running{}{(nothing yet, the tests first)}

First the filter itself, without ROS. In the Octave tutorial this was the loop of
\datei{kalman\_intro\_template.m}; now it is a class with two methods, \thema{predict} and \thema{correct}, and the same equations for every sensor.

\begin{script}{section 3.1, prediction and correction}
Prediction from step \(n-1\) to step \(n\), here without a known input (\(\mat B \hat{\vec u}_n = 0\)):
\begin{equation}
  \tilde{\vec x}_n = \mat A\,\tilde{\vec x}_{n-1}^+ ,
  \qquad
  \tilde{\mat P}_n = \mat A\,\tilde{\mat P}_{n-1}^+\,\mat A\T + \mat Q .
  \tag{3.6, 3.8}
\end{equation}
Correction with the measurement \(\hat{\vec y}_n\) and the predicted measurement
\(\tilde{\vec y}_n = \mat H\,\tilde{\vec x}_n\) (3.7):
\begin{align}
  \mat K_n &= \tilde{\mat P}_n \mat H\T \big(\mat H\,\tilde{\mat P}_n \mat H\T + \mat R\big)^{-1} ,
  \tag{3.10}\\
  \tilde{\vec x}_n^+ &= \tilde{\vec x}_n + \mat K_n \big(\hat{\vec y}_n - \tilde{\vec y}_n\big) ,
  \tag{3.11}\\
  \tilde{\mat P}_n^+ &= \big(\mat I - \mat K_n \mat H\big)\,\tilde{\mat P}_n .
  \tag{3.12}
\end{align}
The constant-velocity model of the worked example, here with a step of \(\Delta t\) seconds:
\begin{equation}
  \begin{pmatrix} x\\ y\\ v_x\\ v_y \end{pmatrix}_n =
  \underbrace{\begin{pmatrix} 1&0&\Delta t&0\\ 0&1&0&\Delta t\\ 0&0&1&0\\ 0&0&0&1 \end{pmatrix}}_{\mat A}
  \begin{pmatrix} x\\ y\\ v_x\\ v_y \end{pmatrix}_{n-1} .
  \tag{3.13}
\end{equation}
\end{script}

\begin{todo}{1.1}
\datei{models.py}, \thema{cv\_model(dt)}: \(\mat A\) of (3.13) as a \(4 \times 4\) numpy array.
Check: \thema{cv\_model(0.5) @ [0, 0, 2, 4]} is \thema{[1, 2, 2, 4]}.
\end{todo}

\begin{todo}{1.2}
\datei{kalman\_filter.py}, \thema{predict(A, Q)}: the state and the covariance with (3.6, 3.8).
\end{todo}

\begin{todo}{1.3}
\thema{correct(z, H, R)}: the gain \(\mat K\) with (3.10). The innovation and its covariance
\(\mat S = \mat H \tilde{\mat P} \mat H\T + \mat R\), the bracket of (3.10), are already computed.
\end{todo}

\begin{todo}{1.4}
\thema{correct}: the new state (3.11) and the new covariance (3.12).
\end{todo}

\(\mat Q\) is given: \thema{process\_noise(dt, sigma\_a)} in \datei{models.py}, more in \cref{sec:task1c}.


\thema{pytest -k task1} also tries two examples of the script: the covariance ellipse (\(\tilde{\mat P} =
\big(\begin{smallmatrix}2.56 & 1.2\\ 1.2 & 1.0\end{smallmatrix}\big)\), only \(x\) measured with
\(\sigma = 0.6\), gives \(\mat K = (0.877,\ 0.411)\T\)), and the worked example (after 100 fixes the velocity is
\((1,\ 0.5)\)). Test 1.5 fails until \cref{sec:task1b}.

\begin{checkpoint}
\thema{pytest -k "task1 and not 5"} passes. Why does a measurement of \(x\) alone change the estimate of
\(v_x\)? Look at \(\tilde{\mat P}\) before the correction.
\end{checkpoint}

\newpage
\section{Task 1: GPS alone, in the loop\minutes{15}}
\label{sec:task1b}
\running{}{-p task:=1}

The node needs \(\hat{\vec y}\), \(\mat H\) and \(\mat R\) of every fix. The GPS measures the position, so
\(\mat H\) picks \(x\) and \(y\) out of the state. \(\mat R\) comes with the message: the simulator puts the
variances of the fix into its \thema{covariance}.

\begin{todo}{1.5}
\datei{models.py}, \thema{gps\_measurement(msg)}: \thema{z}, \thema{H} (\(2 \times 4\)) and \thema{R}
(\(2 \times 2\)) of one fix. The covariance of the message is a \(6 \times 6\) matrix row by row, in the order
\(x, y, z\), roll, pitch, yaw (\cref{tab:sensors}).
\end{todo}

Start your node. It starts at the first fix with the velocity 0 and \(\sigma_v = 1\)\,m/s. Then it predicts
ten times per second and corrects with every fix. Every 5\,s it prints the state, \(\sigma\) and the gain
\(K\) of \(x\).
\begin{itemize}
  \item Watch the red ellipse in RViz: it grows between two fixes and shrinks at every fix. The script says
        \(\mat P\) \emph{breathes}.
  \item How long does it take until the velocity in terminal 2 is about 0.5\,m/s?
  \item \(K\) settles to a fixed value. Which one? What does it mean that it is not 1?
\end{itemize}

\expected[]{task1}

\begin{checkpoint}
\thema{pytest -k task1} passes. The filter is about two times better than the GPS, and the NEES is about 1.
\end{checkpoint}

\newpage
\section{Task 1: GPS alone, the process noise\minutes{30}}
\label{sec:task1c}
\running{}{-p task:=1 -p sigma\_a:=0.0}

\begin{note}
Given in \datei{models.py}: \thema{process\_noise(dt, sigma\_a)}. The model says ``the velocity stays''. It does
not: the robot brakes and turns. \(\mat Q\) says by how much the model may be wrong: an unknown acceleration with
standard deviation \(\sigma_a\), constant over the step, moves the position by \(a\,\Delta t^2/2\) and the velocity
by \(a\,\Delta t\). Per axis
\(\mat Q = \sigma_a^2 \, \vec g \vec g\T\) with \(\vec g = (\Delta t^2/2,\ \Delta t)\T\).
\(\sigma_a = 0\) is \(\mat Q = \mat 0\).
\end{note}

\(\sigma_a\) is the only number you tune in task 1. Run the node with each \(\sigma_a\) of \cref{tab:q}, for a
full lap each (70\,s), and take the last line of the monitor. Ctrl-C, then start again with the next value.

The monitor checks two things. The error: how far your estimate is from the truth. And the honesty: is the
\(\sigma\) your filter claims as large as its error? The NEES (normalized estimation error squared) is
\begin{equation*}
  \text{NEES} = \tfrac12\,\vec e\T \mat P_{xy}^{-1}\,\vec e , \qquad \vec e = \text{estimate} - \text{truth},
\end{equation*}
averaged over the last 30\,s. About 1: \(\sigma\) is honest. Much larger than 1: the filter is sure of itself and
wrong.

\begin{table}[ht]\centering
\renewcommand{\arraystretch}{1.4}
\caption{GPS alone, five values of \(\sigma_a\). Fill in.}
\label{tab:q}
\begin{tabular}{@{}p{2.4cm}p{1.9cm}p{1.9cm}p{1.9cm}p{1.9cm}p{1.9cm}@{}}
\toprule
\(\sigma_a\) in m/s\(^2\) & GPS in m & filter in m & max in m & \(\sigma\) in m & NEES \\
\midrule
0 & & & & & \\
0.1 & & & & & \\
0.3 & & & & & \\
1 & & & & & \\
3 & & & & & \\
\bottomrule
\end{tabular}
\end{table}

\begin{itemize}
  \item \(\sigma_a = 0\) is \(\mat Q = \mat 0\), the first exercise of the script. What happens to \(K\) and
        to \(\sigma\) after a few laps? And to the error at the corners and at the stop?
        \thema{plot\_run.py} shows it in the velocity.
  \item Which \(\sigma_a\) gives the smallest error? Which one gives NEES \(\approx 1\)? Compare with what the
        robot does: it brakes and starts with 0.3\,m/s\(^2\), and in a corner (\(r = 1\)\,m) it accelerates
        with \(v^2/r = 0.25\)\,m/s\(^2\) towards the centre.
  \item \(\sigma_a = 3\): why is the error larger again, although \(\sigma\) is honest?
\end{itemize}

\expected[\(\sigma_a = 0\); ]{q0}

\begin{checkpoint}
Explain with your table: too small a \(\mat Q\) makes the filter confident and wrong, too large a \(\mat Q\) makes
it honest and noisy. Show the plots of \(\sigma_a = 0\) and \(\sigma_a = 0.3\).
\end{checkpoint}

\newpage
\section{Task 2: GPS and odometry\minutes{35}}
\label{sec:task2}
\running{}{-p task:=2}

Now the node also subscribes to \thema{/muster/odom} and corrects with its velocity, 50 times per second. The
state stays the same, and so do \thema{predict} and \thema{correct}: only \(\hat{\vec y}\), \(\mat H\) and
\(\mat R\) are new. This is the second part of the Octave tutorial: more measurements, same filter.

The odometry gives the velocity \((v_x^R, v_y^R)\) in the frame of the robot (\cref{tab:sensors}). The state
holds it in the world frame. Rotate it with the heading \(\theta\) of the odometry, as in lab 1:
\begin{equation*}
  \begin{pmatrix} v_x \\ v_y \end{pmatrix} =
  \begin{pmatrix} \cos\theta & -\sin\theta \\ \sin\theta & \cos\theta \end{pmatrix}
  \begin{pmatrix} v_x^R \\ v_y^R \end{pmatrix} .
\end{equation*}

\begin{todo}{2.1}
\datei{models.py}, \thema{odom\_measurement(msg, sigma\_v)}: \thema{z} the velocity in the world frame,
\thema{H} (\(2 \times 4\)) picks \(v_x\) and \(v_y\) out of the state, \thema{R} \(= \sigma_v^2\,\mat I\).
\(\theta\), \(v_x^R\) and \(v_y^R\) are already read from the message.
\end{todo}

Start the node with \thema{task:=2}. Then try the values of \(\sigma_v\) in \cref{tab:sigmav}, a lap each.

\begin{table}[ht]\centering
\renewcommand{\arraystretch}{1.4}
\caption{GPS and odometry, \(\sigma_a = 0.3\)\,m/s\(^2\), four values of \(\sigma_v\). Fill in.}
\label{tab:sigmav}
\begin{tabular}{@{}p{2.4cm}p{2.4cm}p{2.4cm}p{2.4cm}p{2.4cm}@{}}
\toprule
\(\sigma_v\) in m/s & filter in m & max in m & \(\sigma\) in m & NEES \\
\midrule
0.01 & & & & \\
0.05 & & & & \\
0.2 & & & & \\
0.5 & & & & \\
\bottomrule
\end{tabular}
\end{table}

\begin{itemize}
  \item How much better than GPS alone (\cref{tab:q}, \(\sigma_a = 0.3\))? Compare the velocity in
        \thema{plot\_run.py}.
  \item The speed of the odometry scatters by about 0.001\,m/s (\cref{tab:sensors}). Yet \(\sigma_v = 0.01\) gives a
        large NEES. The odometry thinks the robot is 3\,\% faster than it is: 0.015\,m/s, the same error in
        every message. Why does the filter, which gets 50 such messages per second, become far too sure of
        itself? What does that mean for choosing \(\mat R\)?
\end{itemize}

\expected[\(\sigma_v = 0.05\); ]{task2}

\begin{checkpoint}
\thema{pytest -k task2} passes. Name the \(\sigma_v\) you keep, and why.
\end{checkpoint}

\newpage
\section{Task 2: the tunnel\minutes{25}}
\label{sec:task2b}
\running{ gps:=tunnel}{-p task:=1 \textrm{, then} -p task:=2 -p sigma\_v:=0.2}

Stop terminal 1 and start it with \thema{gps:=tunnel}. The tunnel covers the upper left corner of the loop and
the stop (\cref{fig:field}): about 18\,s without a fix. Every time the robot leaves the tunnel, the monitor prints
the largest error and the largest \(\sigma\) inside.

Run task 1 for two laps, then task 2 for two laps. Fill in \cref{tab:tunnel}.

\begin{table}[ht]\centering
\renewcommand{\arraystretch}{1.4}
\caption{Through the tunnel. Fill in.}
\label{tab:tunnel}
\begin{tabular}{@{}p{6.0cm}p{3.8cm}p{3.8cm}@{}}
\toprule
& task 1, GPS alone & task 2, with odometry \\
\midrule
largest error in the tunnel, lap 1 & & \\
largest error in the tunnel, lap 2 & & \\
largest \(\sigma\) in the tunnel & & \\
\bottomrule
\end{tabular}
\end{table}

\begin{itemize}
  \item In the tunnel the filter only predicts. What does the constant-velocity model assume the robot does
        there? What does the robot really do?
  \item With \thema{plot\_run.py}: how does \(\sigma\) grow in the tunnel (grey band), with and without
        odometry? Is it honest in both cases?
  \item What happens at the first fix after the tunnel, with GPS alone? Why does this one fix move the
        estimate so much more than the fixes before the tunnel?
\end{itemize}

\expected[task 1; ]{tunnel1}
\expected[task 2; ]{tunnel2}

\begin{checkpoint}
Explain the difference in \cref{tab:tunnel} with the two models of the robot: what the filter knows without
a measurement, in task 1 and in task 2.
\end{checkpoint}

\newpage
\section{Extra task: multipath\minutes{10}}
\label{sec:extra}
\running{ gps:=multipath}{-p task:=2 -p sigma\_v:=0.2 -p gate:=9.21}

\thema{gps:=multipath} adds a glass front on the right straight (\cref{fig:field}). It reflects the GPS signal:
every fix there is 2\,m off. The filter believes these fixes and is pulled away. Run task 2 without the gate
first and look at the monitor line \thema{multipath}.

A fix this far off is unlikely. The innovation \(\vec\nu = \hat{\vec y} - \tilde{\vec y}\) has the covariance
\(\mat S\) (the bracket of 3.10). Its squared Mahalanobis distance
\begin{equation*}
  d^2 = \vec\nu\T \mat S^{-1} \vec\nu
\end{equation*}
is larger than 9.21 for only one fix in a hundred, if the filter is right (\(\chi^2\) distribution, two degrees
of freedom). With \thema{-p gate:=9.21} the node drops every fix above that.

\begin{todo}{Z.1}
\datei{kalman\_filter.py}, \thema{mahalanobis2(innovation, S)}: \(d^2\). \thema{correct} calls it only when
the gate is on.
\end{todo}

\begin{itemize}
  \item How many fixes does the gate drop per lap (terminal 2)? How large is the error at the glass front now?
  \item Try the gate with task 1 (GPS alone). Why does it help much less there?
  \item A gate is dangerous with a filter that is sure of itself and wrong. Try \thema{task:=1},
        \thema{sigma\_a:=0}, \thema{gate:=9.21}. What happens, and why does the filter not recover?
\end{itemize}

\expected[with the gate; ]{gate}

\section{Outlook: the heading}
\label{sec:outlook}
Your filter borrows the heading from the odometry and never questions it. If the odometry turned the robot by a
few degrees too much, the filter would rotate every velocity wrong, and the GPS would have to pull it back all the
time. The cure is to estimate the heading as well: \(\vec x = (x, y, \theta, v)\T\), with the odometry as the
input of the model. That model has \(\cos\theta\) and \(\sin\theta\) in it: non-linear. The Extended Kalman filter
(script, section 3.1.4) linearizes it at the current estimate. The equations stay the ones you wrote today.

\end{document}
